%pacotes
%\usepackage[utf8]{inputenc} % desnecessário no XeLaTeX/LuaLaTeX
\usepackage[brazil]{babel}
\usepackage{graphicx}
\usepackage{changepage}
\usepackage{csquotes}
\usepackage[toc]{appendix}
\usepackage{booktabs} %tabelas
\usepackage{indentfirst} %recuo do paragrafo
\RequirePackage{xparse}


%configuração da página
\usepackage{geometry}
\geometry{
    top=2cm,      % distância do topo da página até o início do cabeçalho
    bottom=3cm,
    left=3cm,
    right=2cm,
    headheight=15pt, % altura do cabeçalho (necessário para fancyhdr)
    includehead       % inclui o cabeçalho no cálculo da margem superior
}

%recuo de paragros e espaço entre eles
\usepackage{indentfirst} % indenta primeiro parágrafo após seção
\setlength{\parindent}{1.25cm} % recuo padrão ABNT
\setlength{\parskip}{0pt} % remove espaço extra entre parágrafos

%Espaçamento entre linhas
\usepackage{setspace}
\onehalfspacing  % 1,5 entre linhas

%configuração da paginação
\usepackage{fancyhdr}
\pagestyle{fancy}
\fancyhf{} % limpa cabeçalho e rodapé
\fancyhead[R]{\thepage}% número da página no topo direito
\renewcommand{\headrulewidth}{0pt} % sem linha
\fancyfoot{} % rodapé vazio

%funcionamento de listas e sumario
\usepackage{tocloft} 

\usepackage[hidelinks]{hyperref} % precisa vir por último


\setlength{\cftbeforesecskip}{0pt}  % espaço antes de cada seção
\setlength{\cftsecindent}{0pt}      % remove indentação da seção
\setlength{\cftsubsecindent}{0pt}   % remove indentação da subseção
\setlength{\cftsubsubsecindent}{0pt}% remove indentação da subsubseção

% seções de nível dois (1.1, 1.2) em negrito conforme definição nas seções
\renewcommand{\cftsubsecfont}{\bfseries}
\renewcommand{\cftsubsecpagefont}{\bfseries}

% pontinhos nas listas e no sumário
\renewcommand{\cftdotsep}{2} % espaço entre os pontos
\renewcommand{\cftsecleader}{\cftdotfill{\cftdotsep}}
\renewcommand{\cftsubsecleader}{\cftdotfill{\cftdotsep}}
\renewcommand{\cftsubsubsecleader}{\cftdotfill{\cftdotsep}}

% largura reservada para o número no sumário e nas listas
\setlength{\cftsecnumwidth}{3em}      % seções 1, 2, etc.
\setlength{\cftsubsecnumwidth}{3em}   % subseções 1.1, 1.2, etc.
\setlength{\cftsubsubsecnumwidth}{3em}% subsubseções 1.1.1, 1.1.2
\setlength{\cftparanumwidth}{3em} % seções 1.1.1.1

% configuração dos links do arquivo
\hypersetup{
    colorlinks=true,    % links aparecem coloridos
    linkcolor=black,    % links internos (capítulos, figuras) 
    citecolor=black,    % citações
    urlcolor=black       % links externos (URL)
}
\urlstyle{same}  % usa mesma fonte do texto normal para URLs

% Remove indentação dos níveis do sumário
\setlength{\cftbeforesecskip}{0pt}  % espaço antes de cada seção
\setlength{\cftsecindent}{0pt}      % remove indentação da seção
\setlength{\cftsubsecindent}{0pt}   % remove indentação da subseção
\setlength{\cftsubsubsecindent}{0pt}% remove indentação da subsubseção

% pacotes para figuras
\usepackage{float} %uso do H 
\usepackage{caption}  % para \captionsetup
\usepackage{setspace} % para espaçamento (\singlespacing, \onehalfspacing, etc.)

% Pacote IFTO
\usepackage{modeloIFTO}

%fonte
\usepackage{fontspec}  % Pacote para fontes do sistema
\setmainfont{Arial}    % Define Arial como fonte principal

%biblatex
\usepackage[
backend=biber,
style=abnt
]{biblatex}

% ajuste em tabelas
\captionsetup[table]{labelsep=endash} % Substitui : por travessão nas tabelas

% configuração dos títulos de seção
\makeatletter
% 1 SEÇÃO PRIMÁRIA: Negrito + Caixa Alta + Tamanho Normal (12)
\renewcommand{\section}{\@startsection{section}{1}{\z@}%
    {-3.5ex \@plus -1ex \@minus -.2ex}%
    {2.3ex \@plus.2ex}%
    {\normalfont\normalsize\bfseries\MakeUppercase}}
% 1.1 SEÇÃO SECUNDÁRIA: Negrito + Caixa baixa + Tamanho Normal (12)
\renewcommand{\subsection}{\@startsection{subsection}{2}{\z@}%
    {-3.25ex\@plus -1ex \@minus -.2ex}%
    {1.5ex \@plus .2ex}%
    {\normalfont\bfseries\normalsize}}
% 1.1.1 Seção Terciária: Sem Negrito + Caixa Baixa + Tamanho Normal (12)
\renewcommand{\subsubsection}{\@startsection{subsubsection}{3}{\z@}%
    {-3.25ex\@plus -1ex \@minus -.2ex}%
    {1.5ex \@plus .2ex}%
    {\normalfont\normalsize}}
\makeatother

%bibliografias
\addbibresource{references.bib}

%notas de rodapé
\makeatletter
\renewcommand\@makefntext[1]{% redefine a indentação das linhas seguintes
  \noindent
  \hangindent=0.5em
  \linespread{1}\selectfont
  \textsuperscript{\@thefnmark}\hspace{0em}#1%
}
\makeatother